% PyRATE-TA — Kinetic analysis and user manual.
% Build with:  pdflatex main.tex  (run twice, for the table of contents)
\documentclass[11pt,a4paper]{article}

\usepackage[UKenglish]{isodate} % To fix the date
\usepackage[T1]{fontenc}
\usepackage{tgtermes}

\usepackage[utf8]{inputenc}
\usepackage[margin=2.5cm]{geometry}
\usepackage{amsmath,amssymb}
\usepackage{graphicx}
\usepackage{xcolor}
\usepackage{booktabs}
\usepackage{listings}
\usepackage[hidelinks]{hyperref}

\usepackage{microtype}        % best single improvement
	\emergencystretch=3em         % last-resort stretch to avoid overfull
	\tolerance=800                % default 200; higher = more flexible
	\hbadness=1000                % silence minor warnings

% --- Code listing style ----------------------------------------------------
\definecolor{codegray}{gray}{0.95}
\definecolor{codekw}{rgb}{0.10,0.30,0.65}
\definecolor{codestr}{rgb}{0.20,0.50,0.20}
\definecolor{codecom}{rgb}{0.45,0.45,0.45}
\lstdefinestyle{py}{
  language=Python,
  backgroundcolor=\color{codegray},
  basicstyle=\ttfamily\small,
  keywordstyle=\color{codekw}\bfseries,
  stringstyle=\color{codestr},
  commentstyle=\color{codecom}\itshape,
  showstringspaces=false,
  breaklines=true,
  frame=single,
  rulecolor=\color{codegray},
  columns=fullflexible,
}
\lstset{style=py}

% --- Convenience macros ----------------------------------------------------
\newcommand{\code}[1]{\texttt{#1}}
\newcommand{\dA}{\ensuremath{\Delta A}}
\newcommand{\Kmat}{\ensuremath{\mathbf{K}}}
\newcommand{\Cmat}{\ensuremath{\mathbf{C}}}
\newcommand{\Smat}{\ensuremath{\mathbf{S}}}
\newcommand{\Dmat}{\ensuremath{\mathbf{D}}}

% --- Feature-status marker -------------------------------------------------
% Sections describing a design that is specified but not yet implemented are
% marked, so the manual never reads as if an engine already existed.
\newcommand{\planned}{\textcolor{codekw}{\textbf{[planned]}}\ }

\title{\textbf{PyRATE-TA}\\[0.3em]
  \large Rate Analysis \& Target-model Engine for Transient Absorption\\
  \large Kinetic analysis and user manual}
\author{Dr. Ricardo J. Fernández-Terán}
\date{\today}

\begin{document}
\maketitle
\tableofcontents
\bigskip

\section{Introduction}
PyRATE-TA fits time-resolved spectroscopy data: multi-exponential decays, global
analysis with the amplitudes solved by variable projection, target analysis
with a general rate matrix, and lifetime density analysis (LDA) over a dense
regularised grid --- together with the decay-, evolution- and species-associated
spectra that come out of them.

PyRATE-TA is the analysis half of a two-project pair. Its sibling,
\textbf{PyMORGAN}, owns loading, processing and plotting: the file-format
readers and their registry, the \code{Dataset1D}/\code{Dataset2D} objects, the
aesthetics layer and every standard plotter. The dependency is strictly
one-way --- \code{pyrate} imports \code{pymorgan}, never the reverse --- and
PyRATE-TA reimplements none of it. A format PyRATE-TA needs and PyMORGAN lacks is
contributed upstream, not forked here. Chapter~\ref{sec:boundary} makes that
boundary explicit, because most confusion about where a piece of code belongs
is resolved by it.

This manual documents the analysis model and the programming interface.
Chapter~\ref{sec:installation} covers installation;
chapter~\ref{sec:boundary} the PyMORGAN boundary;
chapter~\ref{sec:models} the kinetic models and the rate-matrix formalism;
chapter~\ref{sec:fitting} the fitting engines, result objects, session files
and reproducibility requirements;
chapter~\ref{sec:gui} the graphical interface;
chapter~\ref{sec:settings} every setting and the configuration file; and
chapter~\ref{sec:extending} how to add a model, an engine or a GUI panel.

\section{Citing PyRATE-TA}\label{sec:citing}
If you publish results obtained with PyRATE-TA, please cite the following. The
same list is printed by \code{pyrate-ta-cite} and logged when the GUI starts.

\begin{enumerate}
  \item The PyRATE-TA / PyMORGAN software paper --- \emph{manuscript in
    preparation}.
  \item R.\ J.\ Fern\'andez-Ter\'an, E.\ Sucre-Rosales, L.\ Echevarria,
    F.\ E.\ Hern\'andez, \emph{A Sweet Introduction to the Mathematical
    Analysis of Time-Resolved Spectra and Complex Kinetic Mechanisms: The
    Chameleon Reaction Revisited}, J.\ Chem.\ Educ.\ \textbf{2022}, 99,
    2327--2337. \href{https://doi.org/10.1021/acs.jchemed.2c00104}%
    {10.1021/acs.jchemed.2c00104}
  \item I.\ H.\ M.\ van Stokkum, D.\ S.\ Larsen, R.\ van Grondelle,
    \emph{Global and target analysis of time-resolved spectra}, Biochim.\
    Biophys.\ Acta Bioenerg.\ \textbf{2004}, 1657, 82--104.
    \href{https://doi.org/10.1016/j.bbabio.2004.04.011}%
    {10.1016/j.bbabio.2004.04.011} --- global and target analysis, and the
    variable-projection formalism (chapter~\ref{sec:fitting}).
  \item M.\ N.\ Berberan-Santos, J.\ M.\ G.\ Martinho, \emph{The integration of
    kinetic rate equations by matrix methods}, J.\ Chem.\ Educ.\ \textbf{1990},
    67, 375. \href{https://doi.org/10.1021/ed067p375}{10.1021/ed067p375} ---
    the eigenvector solution used by the propagator
    (section~\ref{sec:kmat}).
  \item P.\ C.\ Hansen, \emph{Analysis of discrete ill-posed problems by means
    of the L-curve}, SIAM Review \textbf{1992}, 34, 561--580.
    \href{https://doi.org/10.1137/1034115}{10.1137/1034115} ---
    L-curve corner detection for automatic regularisation in LDA
    (section~\ref{sec:lda}).
  \item G.\ H.\ Golub, M.\ Heath, G.\ Wahba, \emph{Generalized
    cross-validation as a method for choosing a good ridge parameter},
    Technometrics \textbf{1979}, 21, 215--223.
    \href{https://doi.org/10.1080/00401706.1979.10489751}{10.1080/00401706.1979.10489751} ---
    GCV alpha selection for LDA (section~\ref{sec:lda}).
  \item C.\ Slavov, H.\ Hartmann, J.\ Wachtveitl, \emph{Implementation and
    Evaluation of Data Analysis Strategies for Time-Resolved Optical
    Spectroscopy}, Anal.\ Chem.\ \textbf{2015}, 87, 2328--2336.
    \href{https://doi.org/10.1021/ac504348h}{10.1021/ac504348h} ---
    the OPTIMUS software and lifetime density analysis (LDA) formalism as
    applied to ultrafast spectroscopy (section~\ref{sec:lda}).
  \item G.\ F.\ Dorlhiac, C.\ Fare, J.\ J.\ van Thor, \emph{PyLDM -- An open
    source package for lifetime density analysis of time-resolved spectroscopic
    data}, PLOS Comput.\ Biol.\ \textbf{2017}, 13, e1005528.
    \href{https://doi.org/10.1371/journal.pcbi.1005528}{10.1371/journal.pcbi.1005528} ---
    open-source lifetime density analysis algorithms.
\end{enumerate}

A routine implementing a published method also logs its own reference through
\code{pyrate_ta.cite} when it runs, so the log of a fit records which formalism
produced it, and each reference appears only once per session.

\section{Installation}\label{sec:installation}

PyRATE-TA requires \textbf{Python~3.12 or later}. It can be installed directly from \href{https://pypi.org/project/pyrate-ta/}{PyPI} or installed from source using \href{https://docs.astral.sh/uv/}{\textbf{uv}}.

\subsection{Quick install via PyPI (Recommended)}\label{sec:install-pypi}

To install the latest official release into your active Python environment:

\begin{lstlisting}[language=bash, style=py]
pip install pyrate-ta
# or with uv:
uv pip install pyrate-ta

# Register bundled fonts into matplotlib:
pymorgan-install-fonts
\end{lstlisting}

\noindent You can also launch the graphical interface directly without prior installation using \code{uvx}:

\begin{lstlisting}[language=bash, style=py]
uvx --from pyrate-ta pyrate-ta-gui
\end{lstlisting}

\subsection{Installation from Source (Development)}

From the repository root:

\begin{lstlisting}[language=bash, style=py]
uv venv                          # create .venv (Python >= 3.11)
uv pip install -e ".[gui,dev]"   # editable install + dependencies
uv run pymorgan-install-fonts    # PyMORGAN's fonts, in *this* environment
\end{lstlisting}

The font step is not optional in practice. PyRATE-TA uses PyMORGAN's font stack and
ships no second installer, but Matplotlib registers fonts per environment: skip
it and every figure falls back to DejaVu Sans, which is why a PyRATE-TA plot can
look unlike a PyMORGAN one while both claim the same style. The GUI checks at
start-up and logs the command if the preferred fonts are missing.

uv automatically uses the project's \code{.venv}: prefix one-off commands with
\code{uv run} to avoid activating it manually.

\subsection{Co-development with PyMORGAN}

PyRATE-TA depends on PyMORGAN. While the two are developed side by side under
\code{Scripts/}, PyMORGAN is picked up from the sibling checkout as an editable
path dependency, so an upstream edit is visible immediately without a
reinstall:

\begin{lstlisting}[style=py]
# pyproject.toml
[tool.uv.sources]
pymorgan = { path = "../PyMORGAN", editable = true }
\end{lstlisting}

Replace that block with a plain version pin once PyMORGAN is published to an
index. The same editable dependency puts \code{pymorgan-install-fonts} on
PyRATE-TA's path, which is how the fonts are registered for this environment.

\begin{table}[h]
  \centering
  \begin{tabular}{lll}
    \toprule
    Package & Minimum version & Purpose \\
    \midrule
    \code{pymorgan}   & 0.6    & Loading, datasets, plotting, aesthetics \\
    \code{numpy}      & 2.1.0  & Array operations \\
    \code{scipy}      & 1.14.1 & Optimisation and matrix exponentials \\
    \code{tomlkit}    & 0.13   & Settings file read/write \\
    \code{PyQt6}      & 6.7.1  & Graphical interface framework \\
    \code{PyQt6\_sip} & 13.8.0 & Qt bindings interface \\
    \bottomrule
  \end{tabular}
  \caption{Core dependencies installed by \code{uv pip install -e .}}
  \label{tab:dependencies}
\end{table}

\subsection{Console scripts}

The editable install puts the following entry points on the path:

\begin{table}[h]
  \centering
  \begin{tabular}{ll}
    \toprule
    Command & Purpose \\
    \midrule
    \code{pyrate-ta-gui}      & Launch the fitting interface \\
    \code{pyrate-ta-edit-gui} & Open \code{main\_window.ui} in Qt Designer \\
    \code{pyrate-ta-settings} & Standalone editor for the analysis settings \\
    \bottomrule
  \end{tabular}
  \caption{Console scripts, mirroring PyMORGAN's. Presentation settings are
    edited with \code{pymorgan-settings}, not here.}
  \label{tab:scripts}
\end{table}

\subsection{Verifying the installation}

\begin{lstlisting}[language=bash, style=py]
uv run python -c "import pyrate_ta; print(pyrate_ta.__version__)"
uv run pytest
uv run python scripts/run_checks.py
\end{lstlisting}

\code{run\_checks.py} exercises the settings round-trip, the lazy public API,
the PyMORGAN boundary and the integrity of the Qt Designer layout; the Qt
checks are skipped, not failed, on a headless machine.
\code{python scripts/run\_changed\_tests.py} runs only the suites affected by
the current diff, and falls back to the full suite whenever a change cannot be
mapped --- an unmapped change is never reported as tested.

\subsection{Versioning}

Every edit to the codebase must bump the version in
\code{src/pyrate_ta/\_\_about\_\_.py}, using \code{python scripts/bump\_version.py}
rather than editing it by hand. The scheme is \code{0.x.yymmdd.devN} (PEP~440):
\code{yymmdd} is the build date and \code{N} the iteration within that day,
restarting at 1 on a new date. A bare letter suffix (\code{0.1.260728d}) is not
valid PEP~440 and breaks the build. \code{--check} exits non-zero when the
recorded version was not produced today, which suits a pre-commit hook.

\section{The PyMORGAN boundary}\label{sec:boundary}

Most questions of the form ``where does this code belong?'' are answered by
Table~\ref{tab:boundary}. Read it before adding a module.

\begin{table}[h]
  \centering
  \begin{tabular}{lll}
    \toprule
    Concern & Owner & Entry point \\
    \midrule
    File formats, loader registry      & PyMORGAN & \code{pm.load\_1D}, \code{pm.register\_loader} \\
    Dataset objects                    & PyMORGAN & \code{pm.Dataset1D}, \code{pm.Dataset2D} \\
    Background, chirp, solvent         & PyMORGAN & \code{Dataset1D.background\_correct}, \dots \\
    Contour / spectra / kinetics plots & PyMORGAN & \code{Dataset1D.plot\_*} \\
    Species-spectra rendering          & PyMORGAN & \code{plot\_species\_spectra} \\
    Aesthetics, colourmaps, unit labels & PyMORGAN & \code{pm.Settings}, \code{pm.apply\_style()} \\
    \midrule
    Kinetic models, rate matrices      & PyRATE-TA & \code{pyrate_ta.models} \\
    Fitting engines, global \& target  & PyRATE-TA & \code{pyrate_ta.fit} \\
    Fit results, uncertainties         & PyRATE-TA & \code{pyrate_ta.results} \\
    Analysis-only plots                & PyRATE-TA & \code{pyrate_ta.plot} \\
    \bottomrule
  \end{tabular}
  \caption{Ownership of each concern across the two packages.}
  \label{tab:boundary}
\end{table}

Three rules follow.

\paragraph{The dependency is one-way.}
\code{import pymorgan} inside PyRATE-TA is expected; \code{import pyrate_ta} inside
PyMORGAN is a design error. PyMORGAN's \code{plot\_species\_spectra}
deliberately takes plain arrays (\code{Sfit}, \code{Taus}, \code{TauErr},
\code{isFixTau}, \code{modelType}) rather than a PyRATE-TA result object, precisely
so the arrow never has to reverse; the adapter that produces those arrays lives
on the PyRATE-TA side. \code{scripts/run\_checks.py} asserts the direction.

\paragraph{No file readers here.}
Loading goes through \code{pm.load\_1D(path, data\_type=\dots)}, which resolves
the format through PyMORGAN's registry and therefore picks up every format
PyMORGAN supports, present and future. A new format is contributed upstream to
\code{pymorgan.oneD.load}, or registered at runtime with
\code{pm.register\_loader} --- never forked into PyRATE-TA.

\paragraph{One aesthetics system.}
Call \code{pm.apply\_style()} and read \code{pm.get\_settings()} for label
style, colourmap, time-axis convention and figure sizes. PyRATE-TA's own
\code{Settings} holds analysis-only fields (chapter~\ref{sec:settings}).

\subsection{A minimal session}

\begin{lstlisting}
import pymorgan as pm
import pyrate_ta as pr

data = pm.load_1D("scan.pdat", data_type="PDAT")  # PyMORGAN loads
data.background_correct(tmin=-20, tmax=-5)        # PyMORGAN processes
fit  = pr.fit_global(data, n_components=3)        # PyRATE-TA fits
data.plot_species_spectra(*fit.as_species_args()) # PyMORGAN plots
\end{lstlisting}

\subsection{Units}\label{sec:units}

Signal data is in mOD, delays are in the dataset's own time unit and the probe
axis in its native spectral unit. Fitted lifetimes are reported in the
dataset's time unit and are never silently converted: a lifetime printed
without its unit, or converted behind the user's back, is a bug.

\subsection{Extraction of traces}

PyMORGAN removed its own \code{extract\_kinetic}/\code{fit\_kinetics} when
PyRATE-TA took over analysis, so there is no upstream extraction helper to call:
that primitive lives here. For a headless fit, slice the dataset directly
(\code{data.Z[:, idx, detector]}) rather than re-deriving pixel-lookup logic;
\code{Dataset1D.plot\_kinetics} returns the extracted data as its third value
when a figure is wanted as well.

\section{Kinetic models}\label{sec:models}

\subsection{The bilinear model}

Every model in PyRATE-TA produces a concentration matrix \Cmat{} of shape
$N_\text{delays} \times N_\text{components}$. The data matrix is then bilinear,

\begin{equation}
  \Dmat \;\approx\; \Cmat\,\Smat^{\mathsf{T}},
  \label{eq:bilinear}
\end{equation}

with $\Smat$ of shape $N_\text{pixels} \times N_\text{components}$ holding the
component spectra. The kinetics enter only through \Cmat, and the spectra only
through \Smat; this separation is what makes global analysis tractable
(section~\ref{sec:varpro}).

\subsection{Model families}

\begin{table}[h]
  \centering
  \begin{tabular}{lll}
    \toprule
    \code{ModelType} & Scheme & Spectra obtained \\
    \midrule
    \code{Parallel}   & independent decays          & DAS (decay-associated) \\
    \code{Sequential} & $A \to B \to C \to \cdots$  & EAS (evolution-associated) \\
    \code{Target}     & general \Kmat, with branching & SAS (species-associated) \\
    \bottomrule
  \end{tabular}
  \caption{Model families and the meaning of the resulting spectra.}
  \label{tab:models}
\end{table}

The distinction is not cosmetic. A parallel model makes no claim about
connectivity, and its DAS are linear combinations of the true species spectra;
a sequential model assumes a single irreversible chain, and its EAS are the
spectra of the successive evolving states; only a target model with an
explicitly justified \Kmat{} yields spectra of actual species. The
\code{modelType} string handed to PyMORGAN's \code{plot\_species\_spectra}
must reflect which of the three was used, because it drives the legend wording
upstream --- mislabelling a DAS as an SAS is a scientific error, not a
formatting one.

\subsection{Parallel and sequential models}

For a parallel model with lifetimes $\tau_i$,

\begin{equation}
  C_{ni} = \exp\!\left(-t_n/\tau_i\right),
\end{equation}

and for a sequential model the same expression is replaced by the analytic
solution of the chain, which for distinct lifetimes is a sum of exponentials
with the well-known amplitude factors
$\prod_{j\neq i}\left(\tau_i^{-1}-\tau_j^{-1}\right)^{-1}$. Both are special
cases of the rate matrix below, and are implemented as such wherever that costs
nothing, so there is one propagator to test rather than three.

\subsection{Rate matrices}\label{sec:kmat}

A target model is defined by its rate matrix \Kmat, with off-diagonal element
$K_{ij}$ the rate from compartment $j$ to compartment $i$ and the diagonal the
negative sum of all rates leaving that compartment. The populations obey

\begin{equation}
  \frac{\mathrm{d}\mathbf{c}(t)}{\mathrm{d}t} = \Kmat\,\mathbf{c}(t),
  \qquad
  \mathbf{c}(0) = \mathbf{c}_0 ,
\end{equation}

so that $\mathbf{c}(t) = \exp(\Kmat t)\,\mathbf{c}_0$. \Kmat{} is built
explicitly and \Cmat{} obtained by eigendecomposition where \Kmat{} is
diagonalisable, and by a matrix-exponential propagator otherwise.

\paragraph{Exactly degenerate lifetimes.} When two rates coincide, \Kmat{} is
defective: the eigenvector matrix is singular and the eigenvalue solution does
not exist, although the kinetics are perfectly well behaved (for
$A \rightarrow B \rightarrow$ with equal rates the closed form is
$c_B = kt\,e^{-kt}$). PyRATE-TA splits the degeneracy by an antisymmetric diagonal
perturbation of relative size $10^{-6}$ and continues. This is not a cosmetic
choice: an optimiser routinely passes through equal lifetimes on its way to the
optimum, and refusing to evaluate there would abort otherwise sound fits. The
perturbation is far below any experimental precision and is logged.

\paragraph{Near-degenerate lifetimes.} When two eigenvalues approach each
other, the eigenvector matrix becomes ill-conditioned. The symptom in a fit is a
pair of amplitudes of opposite sign that grow without bound while the residual
barely improves. PyRATE-TA guards this case explicitly and warns when two fitted
lifetimes come within \code{degeneracy\_warn\_ratio} of each other
(chapter~\ref{sec:settings}); a fit that trips this warning should be re-run
with one component fewer before its lifetimes are believed.

\subsection{Predefined schemes}

Branched schemes are kept in \code{pyrate_ta.models.TARGET\_SCHEMES}, keyed by a
short identifier: equilibria (\code{A\_eq\_B} and its decaying variants),
chains with an equilibrium (\code{A\_eq\_B\_to\_C}, \code{A\_eq\_B\_eq\_C}),
branchings into a common product (\code{branch\_to\_common\_D}) and parallel
chains (\code{two\_chains\_decaying\_ends}). A scheme carries its own species
and rate counts, and these need not agree: \code{A <=> B; A ->; B ->} has two
compartments and four rate constants. A target model therefore takes its size
from the scheme, not from the requested number of components, and the number of
free lifetimes is the scheme's rate count.

\begin{lstlisting}
import pyrate_ta as pr

model = pr.make_model("Target", scheme="A_eq_B_to_C")
C = model.concentrations(t, taus=[2.0, 5.0, 20.0, 200.0], irf_fwhm=0.15)
\end{lstlisting}

\subsection{Writing a scheme}\label{sec:notation}

A target model is written as its reactions, not as a matrix:

\begin{lstlisting}
A  ->  B   : k1          # a step
B <->  C   : k2, k3      # an equilibrium (forward, backward)
C  ->      : k4          # decay to the ground state
init A = 1               # optional initial populations
\end{lstlisting}

\code{pyrate\_ta.scheme\_from\_text} turns this into a scheme, and the rate matrix
is \emph{derived}: every reaction contributes $-k$ to the reactant's diagonal
and $+k$ to the product's row, so the columns cannot fail to balance. Three
consequences are worth stating, because they are what make the notation
general:

\begin{itemize}
  \item \textbf{Species are named.} \code{S1}, \code{ICT}, \code{T1} appear on
    the diagram, in the lifetime table and in the result, so a scheme reads as
    chemistry rather than as indices.
  \item \textbf{Rates are named, so they can be shared.} The same rate name on
    two arrows is \emph{one} fitted parameter used twice --- the natural way to
    impose ``these two channels have the same rate'', which a rate matrix
    cannot express without a hand-written builder.
  \item \textbf{Any topology.} Chains, branches, equilibria, cycles,
    disconnected pools and non-decaying products are all just lines; nothing is
    special-cased, and \code{scheme\_to\_text} regenerates the notation so a
    scheme round-trips.
\end{itemize}

A malformed scheme is refused with the offending line number, never
half-built; \code{check\_scheme\_text} returns that message without raising,
which is what the editor's \emph{Update K graph} button reports.

\subsection{Scheme diagrams}

\code{pyrate_ta.plot\_scheme} draws the compartment diagram of any model, scheme
or fit result: one node per compartment, one arrow per rate, the ground state
as a separate node, and an equilibrium as two opposed arrows. Arrows are
labelled with the \emph{rate constants themselves} --- $k_1$, $k_2$, and a sum
such as $k_1+k_2$ where a compartment branches --- rather than with numbers, so
the diagram tells you which row of the rate table an arrow corresponds to when
setting values and bounds. The mapping is derived, not assumed: feeding the
scheme's $k \rightarrow \Kmat$ builder a unit vector shows exactly where each
constant appears. Pass \code{rate\_labels="value"} for the lifetime of each
step instead.

Every family is laid out \textbf{horizontally}: population flows left to right
and the ground state sits below, so a branched scheme spreads across a wide,
short figure instead of stacking a column that wastes the width and shrinks the
nodes.

The scheme is recovered from whatever the object carries. A fit records its
scheme twice --- as the notation (\code{scheme\_text}) and as a key --- and the
\emph{text} is parsed first, so a scheme written in the editor keeps its species
names and rate symbols on the diagram. This matters because a typed scheme is
stored under the key \code{"custom"}, which is in no registry: consulting the
key alone made drawing one fail outright.

\subsection{Instrument response}\label{sec:irf}

The instrument response is convolved into \Cmat, not into the data. For
exponential models this is done analytically: the convolution of a Gaussian of
width $\sigma$ centred at $t_0$ with a decaying exponential is an exponentially
modified Gaussian,

\begin{equation}
  \left(g \ast e^{-t/\tau}\right)(t)
  = \tfrac{1}{2}\exp\!\left(\frac{\sigma^{2}}{2\tau^{2}} - \frac{t-t_0}{\tau}\right)
    \operatorname{erfc}\!\left(\frac{1}{\sqrt{2}}\left(\frac{\sigma}{\tau} - \frac{t-t_0}{\sigma}\right)\right),
  \label{eq:emg}
\end{equation}

which is both faster and better conditioned than numerical convolution. The
three handling modes are \code{None} (no IRF), \code{Gaussian}
(equation~\ref{eq:emg}) and \code{Clip} (drop the delays below \code{t\_min}
instead of modelling the response).

\subsection{The coherent artefact}

The signal around $t=0$ contains contributions --- cross-phase modulation,
two-photon absorption, stimulated Raman --- that no kinetic model describes.
Two treatments are supported: clipping the early delays, or adding artefact
components (the IRF and its first derivatives) to the model. Clipping changes
which data the fit saw, so the applied \code{t\_min} is part of the result's
reproducibility payload (section~\ref{sec:repro}) and not merely a plotting
choice.

\subsection{Chirp}

Chirp is corrected \emph{upstream}, in PyMORGAN, before the data reach a fit.
PyRATE-TA does not implement a second chirp correction, and a fit of visibly
chirped data is a data-processing mistake rather than a modelling one.

\section{Fitting}\label{sec:fitting}

\subsection{Variable projection}\label{sec:varpro}

In equation~\ref{eq:bilinear} the spectra $\Smat$ enter linearly. For any fixed
set of non-linear parameters (the lifetimes, $t_0$ and the IRF width) the
best-fit spectra are therefore available in closed form,

\begin{equation}
  \Smat^{\mathsf{T}} = \left(\Cmat^{\mathsf{T}}\Cmat\right)^{-1}\Cmat^{\mathsf{T}}\Dmat ,
  \label{eq:varpro}
\end{equation}

and only the non-linear parameters need to go to the optimiser. This is
\emph{variable projection}: for a 500-pixel, three-component fit it reduces the
parameter vector from about 1500 entries to three, and it conditions the
problem far better than fitting amplitudes and lifetimes jointly. Amplitudes
must not be placed in the optimiser's parameter vector.

In practice equation~\ref{eq:varpro} is evaluated through a least-squares
solver rather than by forming the normal equations, which squares the condition
number of \Cmat.

\subsection{Engines}

\begin{table}[h]
  \centering
  \begin{tabular}{lll}
    \toprule
    Engine & Data & Output \\
    \midrule
    Single-trace fit & one kinetic trace       & lifetimes, amplitudes \\
    Global analysis  & full $\Dmat$, free \Kmat & lifetimes, DAS or EAS \\
    Target analysis  & full $\Dmat$, fixed connectivity & rates, SAS \\
    LDA              & full $\Dmat$, dense grid & lifetime density map \\
    \bottomrule
  \end{tabular}
  \caption{The four fitting engines. Global, target and single-trace share the
    variable-projection driver; LDA is a separate regularised linear solve.}
  \label{tab:engines}
\end{table}

\begin{lstlisting}
# Global analysis
fit = pr.fit_global(data, n_components=3, model_type="Sequential")
print(fit.taus, fit.tau_err)
data.plot_species_spectra(*fit.as_species_args())

# Preview (no optimisation)
prev = pr.preview_global(data, taus=[1.0, 10.0, 100.0])

# Target analysis
fit = pr.fit_target(data, scheme="A_eq_B_to_C", taus=[2.0, 5.0, 20.0, 200.0])

# Single-trace
fit = pr.fit_kinetics(t, y, n_components=2)

# Lifetime density analysis
lda = pr.fit.lda.solve_lda(data, n_taus=100, alpha="auto", non_negative=True)
\end{lstlisting}

Every fit routine is callable headlessly: a fit must be reachable and testable
from a plain script with no \code{QApplication}. The GUI is a caller, never a
prerequisite. Where a long fit puts up a progress dialog, the optimiser reports
through a plain callback and the GUI adapts it; the engine itself does not
touch widgets, and guards on \code{QApplication.instance()} being \code{None}.

\paragraph{Preview mode.} \code{pr.preview\_global} evaluates a model at
the given lifetimes \emph{without} optimising. The spectra still follow from
the data (they are the least-squares solution for the given \Cmat), so a
preview shows what those lifetimes actually imply and what the residual would
be. Nothing is fitted, so there is no convergence report and no uncertainties
--- and the result says so rather than presenting a guess as an answer.

\paragraph{Component ordering.} Variable projection sees only the column
space of \Cmat, so permuting the lifetimes fits equally well and the
optimiser's ordering depends on the initial guess. For parallel and sequential
models the components are sorted by lifetime on output, the spectra permuted to
match, and the fact recorded on the result (\code{sorted\_by\_lifetime}). A
target model is left exactly as fitted: its compartments are defined by the
scheme, and reordering them would misrepresent it.

\subsection{The coherent artefact}\label{sec:artefact}

Around time zero a transient carries signal no kinetic model describes:
cross-phase modulation and two-photon absorption in the solvent and the cell
windows. Left alone, the fit absorbs it into the shortest lifetime, which is
then not a lifetime at all.

Ticking \emph{Fit coherent artefact} (beside the IRF table, or
\code{coherent\_artifact=True} on any engine) appends three columns to the
design matrix: the instrument response and its first two derivatives,
\[
  g(t),\qquad g'(t),\qquad g''(t),
\]
each scaled to unit maximum so the second derivative --- larger than the
Gaussian by $1/\sigma^{2}$ --- does not dominate the conditioning. Their
amplitudes are solved per probe pixel exactly like a species spectrum, so the
artefact gets spectra of its own instead of distorting $\tau_1$. The
description follows Kovalenko \emph{et al.}, Phys.\ Rev.\ A \textbf{59}, 2369
(1999), which the log cites when the basis is first built.

The artefact columns are \emph{not} species: the result records how many there
are (\code{n\_artifact}), \code{as\_species\_args()} hands over only the kinetic
ones --- a column with no lifetime must not be labelled with someone else's ---
and \code{artefact\_spectra} exposes the rest. They are worth a look: an
artefact spectrum with amplitude far from time zero is a sign it is absorbing
real kinetics. The option needs a Gaussian IRF, since the basis is built from
it; without one the checkbox is disabled and the engines refuse rather than
inventing a width.

\subsection{Constrained amplitudes and sign masks}

When physical constraints are imposed on the species spectra --- non-negative
absorption for real excited states, non-positive for a ground-state bleach
contribution --- variable projection solves a bounded linear least-squares
problem per probe pixel rather than an unconstrained linear solve,

\begin{equation}
  \min_{\mathbf{S}_p}\,
    \lVert \mathbf{C}_w\,\mathbf{S}_p - \mathbf{d}_{w,p}\rVert^2
  \quad\text{subject to}\quad
    \mathbf{l} \le \mathbf{S}_p \le \mathbf{u}.
\end{equation}

A \code{sign\_mask} array of length $N_c$ controls this: $+1$ enforces
$S_p \ge 0$, $-1$ enforces $S_p \le 0$, and $0$ is unconstrained. When all
entries are zero the problem collapses to the standard unconstrained solve.
Simple non-negative least squares uses \code{scipy.optimise.nnls}; a mixed
mask uses \code{scipy.optimise.lsq\_linear}. The constrained solver is in
\code{pyrate_ta.fit.constrained}.

\subsection{Ground-state bleach and absolute species spectra}\label{sec:groundstate}

Transient absorption and vibrational data are difference spectra:
$\mathbf{S}_\text{diff} = \mathbf{S}_\text{excited} - f\,\mathbf{GS}$.
Adding back a scaled ground-state spectrum converts them to absolute absorption
spectra:
\begin{equation}
  \mathbf{S}_\text{abs}[:,\,i]
    = \mathbf{S}_\text{diff}[:,\,i] + f_i\cdot a\cdot \mathbf{GS},
\end{equation}
where $a$ is the ground-state scale factor and $f_i$ is the fraction of
ground state bleached by species $i$.

\code{pyrate_ta.fit.groundstate.compute\_absolute\_spectra} performs this
conversion. The maximum physically allowable scale factor $a_\text{max}$ ---
the largest $a$ before any pixel in any absolute spectrum turns negative --- is
computed first, so the automatic choice (\code{gs\_scale="auto"}) selects the
tightest non-negative solution. The result stores \code{S\_abs},
\code{gs\_scale} and \code{a\_max} for plotting and reporting.

\subsection{Lifetime density analysis}\label{sec:lda}

Instead of fitting $N$ discrete lifetimes, LDA solves for a continuous
distribution of amplitudes over a dense, logarithmically spaced grid of fixed
lifetimes,

\begin{equation}
  \min_{\Smat^{\mathsf{T}}}
  \left\lVert \mathbf{W} \odot \left(\Dmat - \Cmat(\boldsymbol{\tau}_\text{grid})\,
  \Smat^{\mathsf{T}}\right)\right\rVert^2
  + \alpha^2 \left\lVert \mathbf{L}\,\Smat^{\mathsf{T}}\right\rVert^2,
  \label{eq:lda}
\end{equation}

where $\mathbf{L}$ is a regularisation operator and $\alpha$ controls the
balance between fit quality and smoothness.

\paragraph{Regularisation operators.} Three are available: the identity
(ridge, $\mathbf{L} = \mathbf{I}$), a first-difference operator ($\mathbf{L} =
\mathbf{D}_1$, penalising slope), and a second-difference operator ($\mathbf{L}
= \mathbf{D}_2$, penalising curvature; the default). The augmented system
$[\Cmat;\,\alpha\mathbf{L}]\,\Smat^{\mathsf{T}} = [\Dmat;\,\mathbf{0}]$ is
solved in one least-squares call.

\paragraph{Automatic alpha selection.} When \code{alpha="auto"} three
methods are available:
\begin{description}
  \item[\code{lcurve}] The corner of the log-residual--log-norm curve ---
    maximum curvature by the Menger formula (Hansen 1992). This is the default.
  \item[\code{gcv}] Generalised cross-validation (Golub et al., 1979).
  \item[\code{morozov}] Morozov discrepancy principle: match the residual norm
    to the estimated noise level, estimated from the pre-time-zero region.
\end{description}

\paragraph{Additional options.}
\begin{description}
  \item[\code{non\_negative}] Enforce $S \ge 0$ per probe pixel via NNLS
    after the alpha scan. Physically motivated for some spectral regions.
  \item[\code{svd\_components}] Pre-filter the data matrix with the top $K$
    SVD components before constructing \Cmat, suppressing noise-dominated
    singular values.
  \item[\code{n\_bootstraps}] Monte Carlo resampling of the residuals:
    refits the LDA \code{n\_bootstraps} times and records the standard
    deviation of the integrated amplitude $A(\tau) = \sum_p |S_{p,i}|$ as a
    confidence estimate.
  \item[\code{find\_peaks}] Automatic peak centroid detection from
    $A(\tau)$, using \code{scipy.signal.find\_peaks} with a 5\% prominence
    threshold.
  \item[\code{coherent\_artifact}] Appends the same IRF-derivative basis as
    the discrete engines (section~\ref{sec:artefact}) to absorb time-zero
    contributions, leaving the regularised columns for the kinetics only.
\end{description}

The result is an \code{LDAResult} object (section~\ref{sec:results-lda})
carrying the density map, the alpha scan diagnostics, residuals, and optional
bootstrap errors. It exports the map to PyMORGAN's PDAT format via
\code{.to\_pdat(path)}.

\subsection{Watching a fit run}\label{sec:monitor}

A global fit blocks for seconds to minutes, and the cost alone says little
about what is going wrong. Setting \code{[gui] fit\_monitor\_every = N} opens a
window showing one bar per parameter, its current value above it and its name
below, refreshed every $N$ residual evaluations; $N = 0$ (the default) never
opens it. The scale is symlog by default
(\code{fit\_monitor\_scale}), which is the only one that shows a lifetime of
3000 and a time zero of $\approx 0$ on the same axis. Fixed parameters are drawn
greyed, so a bar that never moves is not mistaken for one that has converged.

The rate is a setting because each update costs a redraw \emph{and} an
event-loop pump: at $N = 1$ a fast fit spends more time drawing than fitting,
while $N = 5$--$10$ is enough to watch a lifetime run away to the edge of the
delay window. The drawing lives in \code{pyrate_ta.plot.monitor} and is Qt-free, so
the same bars --- and \code{plot\_parameter\_history}, the trajectory against
evaluation number --- can be produced from a script.

\subsection{Weighting and the fit statistic}\label{sec:weights}

If the dataset carries a per-point noise estimate --- a sibling \code{.pdatn}
file loaded alongside the \code{.pdat}, or the spread over single scans --- the
residuals can be weighted by it. PyMORGAN supplies the array through
\code{Dataset1D.noise\_array()}; PyRATE-TA never estimates noise itself.

With weights $w = 1/\sigma$, and $N$ the number of points actually fitted and
$p$ the number of free parameters, the statistic is a true reduced
chi-squared,

\begin{equation}
  \chi^{2}_{\nu} = \frac{1}{N-p}\sum_{n,k}
    \left(\frac{D_{nk} - \left(\Cmat\Smat^{\mathsf{T}}\right)_{nk}}{\sigma_{nk}}\right)^{2},
  \label{eq:chi2red}
\end{equation}

so that $\chi^{2}_{\nu} \approx 1$ means the residuals sit at the noise level,
$\gg 1$ that the model is too simple (or the noise underestimated), and
$\ll 1$ that the data are being overfitted (or the noise overestimated).
Without weights the reported statistic is the plain sum of squared residuals,
which carries no such interpretation and is not comparable between datasets.
The read-out in the GUI is therefore labelled \emph{SSR} or \emph{red.\
chi2} according to which one is in force, and never simply ``$\chi^2$''.

The parameter count $p$ includes the $N_p N_c$ linear amplitudes: variable
projection solves them in closed form, but it still spends those degrees of
freedom, so counting only the non-linear parameters would flatter the fit. The
value obtained with the non-linear-only denominator is reported alongside
(\code{FitStatistic.value\_nonlinear\_dof}), because the two differ by a large
factor and the choice should never be implicit.

Two practical points. Noise values are clipped from below at
\code{noise\_floor\_fraction} of the median, because a single near-zero
$\sigma$ would otherwise dominate the cost function entirely; and points with
non-finite or non-positive $\sigma$ are dropped from the fit along with their
data, with the count recorded on the result. Weights are diagonal: correlation
of the noise between pixels is ignored, so $\chi^{2}_{\nu}$ should not be read
as more rigorous than that assumption allows.

Requesting weights for a dataset that has no noise array raises. Falling back
silently to an unweighted fit while still calling the result a reduced
chi-squared would be a misreported statistic, which is worse than a refusal.

\subsection{The fixed/free mask}

A caller that knows only about the lifetimes --- the GUI reads its mask from the
lifetime table --- may pass a mask of ``n\_lifetimes'' entries while the
parameter vector is longer. Time zero and the IRF width are appended to that
vector \emph{only when they are being fitted}, so the entries that follow are
free by construction and the mask is padded rather than refused. A mask of any
\emph{other} length is still an error, naming the parameters expected: padding
must complete an incomplete answer, not reinterpret a wrong one.

A \code{fixed} mask may also be a sequence of parameter \emph{names}
(e.g.\ \code{["tau1", "t0"]}), in which case the driver matches by name and
the rest are free.

\subsection{Uncertainties}

One-sigma uncertainties are reported from the linearised covariance estimate at
the optimum. They are conditional on the model being correct, and they are not
confidence intervals for the choice of model: a three-component fit of
two-component data can return small uncertainties on three meaningless
lifetimes. Parameters held fixed receive no uncertainty and are flagged as
fixed, so that a fixed lifetime can never be mistaken for a well-determined
one.

\subsection{Convergence}\label{sec:convergence}

A non-converged fit raises, or is flagged on the result object; it is never
returned as though it had converged. The optimiser's own termination report
(the reason for stopping, the iteration count, the final cost) travels with the
result. Hitting \code{max\_iterations} is a non-convergence, not a result.

Parameters resting on a finite bound at termination are reported separately
(\code{FitReport.at\_bounds}); a bounded optimum is a constraint, not an
optimum, and downstream code should treat it as such.

\subsection{Results}\label{sec:results}

Three result classes share the common fields of \code{KineticFit}:

\begin{description}
  \item[\code{KineticFit}] Single-trace or multi-trace result: lifetimes,
    amplitudes, residuals, reproducibility payload.
  \item[\code{GlobalFit}] Adds nothing structurally (the spectra live in
    \code{S}), but is typed distinctly so a function can assert it received
    a global result rather than a single-trace one.
  \item[\code{TargetFit}] Adds \code{scheme\_key}, \code{scheme\_label},
    \code{scheme\_text}, the rate matrix \code{K} and \code{eigenvalues}.
    The eigenvalues are stored because they, not the rates, are what the data
    actually determine for a branched scheme.
\end{description}

Every result exposes \code{as\_species\_args()} (the adapter for PyMORGAN's
\code{plot\_species\_spectra}), \code{spectra\_kind} (DAS / EAS / SAS),
\code{artefact\_spectra} (the coherent-artefact columns, if any),
\code{summary()} and \code{log\_summary()}.

\subsubsection{LDA result}\label{sec:results-lda}

\code{LDAResult} carries the density map \code{S\_map} of shape
$(N_p \times M)$, the lifetime grid \code{tau\_grid} of length $M$, the chosen
$\alpha$, the L-curve diagnostics, optional bootstrap errors, and a list of
detected peaks. It also carries \code{citations()} and \code{format\_citations()},
which enumerate the references for whichever algorithm combination was used, and
\code{to\_pdat(path)}, which exports the 2D map to PyMORGAN's PDAT format for
further plotting.

\subsection{Reproducibility}\label{sec:repro}

A result object carries enough to reproduce itself:

\begin{itemize}
  \item the model identity (family, number of components, the \Kmat{} where applicable);
  \item the initial guesses and the bounds;
  \item which parameters were held fixed;
  \item the delay and probe ranges actually fitted, including any clipping applied for the coherent artefact;
  \item the solver's convergence report.
\end{itemize}

A lifetime reported without its uncertainty and its fixed/free flag is not a
result. Results serialise through \code{pyrate_ta.io}, which handles fit sessions
--- parameters, masks, results --- and not raw spectroscopy formats.

\subsection{Fit session files}\label{sec:sessions}

\code{pr.save\_fit(path, fit)} writes a \code{.prfit} file: a compressed
\code{.npz} archive holding the arrays (spectra, concentration profiles,
residuals, axes) and the reproducibility payload as a JSON string inside it.
Keeping them together means a saved fit cannot lose its provenance.

\code{pr.load\_fit(path)} reopens a session as a \code{LoadedFit}: the arrays
as attributes, the JSON payload as \code{.meta}. It is deliberately \emph{not}
a live \code{KineticFit} --- a reopened session has no dataset behind it and
cannot be refitted as it stands, and returning something that looks like a
fresh result would invite exactly that. Everything needed to set the fit up
again is in \code{.meta}. A target fit reopens with its \code{scheme\_text},
so the scheme is editable without re-entering it.

\begin{lstlisting}
path = pr.save_fit("scan_fit", fit)     # -> scan_fit.prfit
prev = pr.load_fit("scan_fit.prfit")
print(prev.summary())                    # same format as a live result
print(prev.taus)                         # numpy arrays as attributes
\end{lstlisting}

The format version is embedded in every file and checked on load: a file
written by a newer PyRATE-TA that the current build cannot interpret raises rather
than silently mis-reading it.

\subsection{Testing}

Synthetic-recovery tests are the backbone of the test suite: build data from
known parameters, fit it, and assert that the parameters come back within
tolerance. A fitting project without these is untested regardless of its line
coverage. The failure modes are tested too --- a non-converged fit must be
flagged rather than silently returned, and near-degenerate lifetimes must not
produce exploding amplitudes without a warning.

\section{The graphical interface}\label{sec:gui}

Launch with \code{uv run pyrate-ta-gui} (or \code{python -m pyrate_ta.gui}). The
window reuses
PyMORGAN's plot widget, plot-controls panel and theme so the two applications
look identical side by side.

\subsection{Layout}

The window is a single narrow control column on the left and a plot area on the
right; nothing is scattered across the width, so the default geometry stays
compact.

\begin{description}
  \item[Left column] Four stacked groups: \emph{Data} (load, clear and two
    lamps --- one for the dataset, one for the per-point noise; there is no
    format selector, since a fit needs a PDAT and its \code{.pdatn} sibling is
    picked up automatically), \emph{Model Selection}
    (\emph{Parallel} / \emph{Sequential} / \emph{Target} / \emph{ODE-defined}),
    \emph{Components} (the lifetime table with its bounds and fixed flags, and
    the add/remove buttons), \emph{IRF and time-zero}, and \emph{Fit options}
    (algorithm, restrict-to-display, noise weighting). The algorithm defaults
    to Trust Region Reflective: lifetimes are bounded below by zero, and
    MINPACK's Levenberg-Marquardt cannot honour bounds at all, so choosing it
    with finite bounds is refused rather than silently ignoring them.
  \item[Fit control] At the foot of the left column: the data / fit / residuals
    view selector, \emph{Preview}, \emph{FIT DATA}, \emph{Reset}, the statistic
    read-out, \emph{Copy taus}, and the two post-fit figures ---
    \emph{Plot DAS} (named after the model: DAS, EAS or SAS) and
    \emph{Plot Conc. Profile}. Both open in their own window, and both are
    opened automatically when a fit finishes; each has its own switch
    (\code{plot\_species\_after\_fit}, \code{plot\_profiles\_after\_fit}) in the
    \code{[gui]} section of \code{settings.toml}. A preview does not trigger
    them: it is a look at a guess, not a result.
  \item[Plot area] One figure, not three, so the window carries a single
    navigation toolbar. Its axes form a $3\times3$ grid: the contour takes the
    top-left $2\times2$ block (twice the size of the others), the kinetics
    panel the $2\times1$ column beside it and the transient spectra the
    $1\times2$ row below. The bottom-right cell holds no axis --- the
    \emph{Additional Plots and cuts} group sits there.

    The kinetics panel is drawn \emph{rotated} (PyMORGAN's
    \code{swap\_axes}), so its delay axis runs vertically like the contour's.
    Both shared axes then line up --- probe between the contour and the
    spectra, delay between the contour and the kinetics --- and the duplicated
    tick labels are dropped. Measured data are drawn as translucent points and
    the fit as a thin line over them, both configurable in the \code{[plots]}
    section of \code{settings.toml}.

    The legends of these panels are always drawn \emph{inside} their axes,
    whatever \code{[plot] legend\_location} says in PyMORGAN's settings: three
    panels share one figure, so a legend parked beside an axis lands on its
    neighbour or off the canvas. The setting still governs every stand-alone
    figure PyRATE-TA opens --- pop-out cuts, species spectra, concentration
    profiles --- where beside the axis is PyMORGAN's default and usual look.
  \item[Bottom] The plot-controls panel: axis limits, colour scale, contour
    levels, time-axis convention and probe units.
\end{description}

The embedded figure sets its margins explicitly (\code{[plots] panel\_left},
\code{panel\_right}, \code{panel\_top}, \code{panel\_bottom} and the two gaps)
rather than calling \code{tight\_layout}: it shares a canvas with the controls
panel, and the automatic layout leaves the axes noticeably smaller than the
window allows.

The four \emph{Additional plots and cuts} buttons open stand-alone PyMORGAN
figures. The two cut plots ask which cuts to draw, pre-filled with the
crosshair's position: the crosshair picks one cut, which is what the embedded
panels show, whereas a separate figure is where several are compared. The
prompt is PyMORGAN's (\code{gui.dialogs.ask\_values}), so the syntax ---
numbers, ranges such as \code{1900:5:1950}, or \code{all} --- is the same as in
PyMORGAN's own cut plots.

\subsection{Fitting on a small screen}

The window opens at $940 \times 900$ and is clipped at start-up to the
\emph{available} screen area --- which already excludes the menu bar and the
dock --- then centred, so a geometry saved on a large monitor cannot reopen
half under the dock on a laptop.

Being able to \emph{ask} for a small window is not enough: a layout whose
minimum size hint exceeds the screen simply refuses to shrink. The left column
therefore lives in a scroll area, so the window's minimum height is set by the
plot canvas alone and the controls stay reachable by scrolling when there is not
room for all of them at once.

Buttons are coloured by function, following the same palette as PyMORGAN --- a
contour button looks like a contour button in both applications --- and the
mapping lives in \code{gui/mw\_common.py} rather than in the \code{.ui}, so it
stays theme-responsive. Controls whose feature is not implemented yet are
disabled and carry the reason in their tooltip, rather than being present and
silently doing nothing.

\subsection{The K-matrix editor}

\emph{Define model...} opens the scheme editor: the scheme as text on the left
(section~\ref{sec:notation}), its graph and the derived rate matrix on the
right. \emph{Update K graph} redraws \emph{and} checks --- a scheme that does
not parse produces its error, with the line number, instead of a picture, so
the button answers ``would this scheme work?'' before a fit is started, and
\emph{OK} stays disabled until it does. The matrix is shown with every rate set
to $1$, since what is being checked there is the structure, not the values.
Accepting the scheme sets the lifetime table to one row per rate constant,
named after it.

\emph{Show kinetic graph}, beside it, draws the scheme of whatever model is
currently selected --- compartments, arrows and rate \emph{symbols}
($k_1$, $k_2$, \dots, not their values) --- in its own window. It works before
any fit, from the lifetimes as typed, so a scheme can be checked while it is
being built; after a fit it draws the fitted one.

\subsection{The LDA tab}

A separate \emph{LDA} tab gives access to Lifetime Density Analysis
(section~\ref{sec:lda}) without switching away from the same dataset. Its
controls map directly to the keyword arguments of \code{solve\_lda}:

\begin{description}
  \item[Grid] Number of log-spaced lifetime points and the
    $[\tau_\text{min},\,\tau_\text{max}]$ range (auto-filled from the delay axis).
  \item[Regularisation] Penalty type (\code{ridge} / \code{d1} / \code{d2}) and
    the alpha method (\code{lcurve} / \code{gcv} / \code{morozov} / \code{fixed}).
    When \emph{fixed} is selected an alpha entry box appears.
  \item[Options] Non-negative constraint, SVD pre-filtering (number of
    components), coherent-artefact suppression, and bootstrap iterations.
  \item[Run LDA] Executes \code{solve\_lda} headlessly and populates the
    results table and the lifetime-density map.
\end{description}

After a successful run the tab shows: the integrated amplitude spectrum
$A(\tau) = \sum_p |S_{p}(\tau)|$, the L-curve or GCV score plot (for diagnosing
the alpha choice), and a 2D heatmap $S(\tau,\lambda)$. Detected peaks are
listed in a table with their centroid lifetimes. The result can be exported to a
\code{.pdat} file with \emph{Export map\ldots} for further plotting in PyMORGAN.

\subsection{Reading a fit}

Two figures carry the result, and both are opened automatically when the fit
converges (buttons reopen them at will):

\begin{description}
  \item[Species spectra] The amplitude spectrum of each component, drawn by
    PyMORGAN's \code{plot\_species\_spectra} through the result's
    \code{as\_species\_args()}. Their name follows the model, and so does the
    button: DAS for parallel, EAS for sequential, SAS for target. The
    distinction is not cosmetic --- it is what the spectra may be interpreted
    as.
  \item[Concentration profiles] The populations $C(t)$ of the model itself: one
    of the few views PyMORGAN has no concept of, since they are never measured.
    Component $i$ carries the same colour as its spectrum, and the delay axis
    comes from PyMORGAN's \code{apply\_delay\_axis}, so the profiles can be laid
    beside the contour and read on the same time axis. For a sequential or
    target model this is where the plot earns its keep: it shows \emph{when}
    each species is present, and therefore which spectrum can be trusted where.
\end{description}

\subsection{Layout is declared, not built}

All widgets, layouts, containers and controls are declared in the Qt Designer
file \code{src/pyrate_ta/gui/main\_window.ui}. Python code binds signals, slots,
dynamic styles and data states to widgets defined there; it does not construct
them. Edit the layout with

\begin{lstlisting}[language=bash, style=py]
uv run pyrate-ta-edit-gui        # or: python -m pyrate_ta.gui.designer
\end{lstlisting}

The plot area is promoted to PyMORGAN's \code{WidgetPlot}, and the whole
plot-controls group (\code{PC\_box}) is the one from PyMORGAN's layout, name for
name, so that \code{pymorgan.gui.plot\_controls.PlotControlsPanel} drives it
unmodified. Renaming a \code{PC\_*} widget breaks that contract; a test asserts
the full set is present.

\subsection{Structure of the code}

\code{MainWindow} is assembled from per-tab mixins in \code{gui/tabs/}:

\begin{description}
  \item[\code{DataTabMixin}] Loading, dataset display, the embedded
    data / fit / residuals view and the additional cut plots.
  \item[\code{ModelTabMixin}] Model selection, the lifetime table, bounds,
    fixed flags, IRF and time-zero controls.
  \item[\code{FittingTabMixin}] Fit control: run, preview, reset, statistic
    read-out, copy taus, auto-plots.
  \item[\code{LDATabMixin}] The LDA controls and result display.
\end{description}

A mixin is a plain method container with no state of its own, so \code{self} is
always the main window. \code{main\_window.py} keeps construction, wiring,
menus and theming only; shared constants live in \code{gui/mw\_common.py}.

Wiring is defensive: a connection is made only if the named widget exists, so
the window still loads while the layout is being edited in Designer, and the
missing widget is logged rather than raising at start-up.

\subsection{Rendering}

A contour re-plot cannot be updated artist by artist --- the level set itself
changes --- so rapid control changes, such as a colour-scale slider drag that
emits one signal per step, are coalesced into a single render through a
short-interval timer. Limit changes that do not alter the levels re-apply the
axis limits without re-plotting.

\subsection{Start-up cost}

\code{gui/app.py} imports neither \code{pyrate} nor \code{pymorgan} at module
level, so the window (and, where enabled, the splash screen) can appear before
scipy, matplotlib and the loader registries are paid for. The same discipline
applies to the package itself: \code{pyrate/\_\_init\_\_.py} binds its exports
through a PEP~562 \code{\_\_getattr\_\_}, so \code{import pyrate_ta} is nearly
free and each module is imported on first use of a name it provides.

\section{Settings}\label{sec:settings}

\subsection{Two settings objects}

Presentation --- label style, colourmap, time-axis convention, figure sizes ---
belongs to \code{pymorgan.Settings} and is read through
\code{pm.get\_settings()} after \code{pm.apply\_style()}. PyRATE-TA's own
\code{Settings} holds \emph{analysis-only} fields, listed in
Table~\ref{tab:settings}. There is deliberately no second aesthetics system.

\begin{table}[h]
  \centering
  \begin{tabular}{llll}
    \toprule
    Field & Default & Group & Meaning \\
    \midrule
    \code{n\_components}        & 2            & fit    & Components proposed for a new fit \\
    \code{model\_type}          & Sequential   & fit    & Default model family \\
    \code{irf\_mode}            & Gaussian     & fit    & \code{None} / \code{Gaussian} / \code{Clip} \\
    \code{irf\_fwhm}            & \code{None}  & fit    & Fixed IRF width; \code{None} fits it \\
    \code{t\_min}               & \code{None}  & fit    & Clip delays below this (mode \code{Clip}) \\
    \code{ftol}                 & $10^{-8}$    & solver & Termination tolerance on the cost \\
    \code{xtol}                 & $10^{-8}$    & solver & Termination tolerance on the parameters \\
    \code{max\_iterations}      & 2000         & solver & Cap on iterations; exceeding it is a failure \\
    \code{variable\_projection} & \code{True}  & solver & Solve amplitudes in closed form \\
    \code{report\_uncertainties}& \code{True}  & solver & Report linearised 1-$\sigma$ errors \\
    \code{use\_noise\_weights}  & \code{False} & solver & Weight residuals by the per-point noise \\
    \code{noise\_floor\_fraction} & $10^{-3}$  & solver & Clip $\sigma$ below this fraction of the median \\
    \code{degeneracy\_warn\_ratio} & 1.2       & solver & Warn when two lifetimes come this close \\
    \midrule
    \code{default\_datadir}     & \code{None}  & paths  & Directory offered by the file dialog \\
    \bottomrule
  \end{tabular}
  \caption{The analysis settings that appear in the settings dialog. LDA
    parameters are not listed here; they are set via the GUI LDA tab or
    directly in \code{settings.toml} for headless use
    (Table~\ref{tab:lda-settings}).}
  \label{tab:settings}
\end{table}

\begin{table}[h]
  \centering
  \begin{tabular}{llll}
    \toprule
    Field & Default & Meaning \\
    \midrule
    \code{lda\_n\_lifetimes}      & 100          & Grid points in the lifetime axis \\
    \code{lda\_tau\_min}          & \code{None}  & Grid lower bound; \code{None} = auto \\
    \code{lda\_tau\_max}          & \code{None}  & Grid upper bound; \code{None} = auto \\
    \code{lda\_regularisation}   & L-curve      & Alpha selection: \code{L-curve} / \code{GCV} / \code{Fixed} \\
    \code{lda\_alpha}            & $10^{-3}$    & Strength used when \code{lda\_regularisation} is \code{Fixed} \\
    \code{lda\_non\_negative}    & \code{False} & Constrain distribution to $\ge 0$ \\
    \bottomrule
  \end{tabular}
  \caption{LDA default settings, stored in the \code{[lda]} section of
    \code{settings.toml}. These do \emph{not} appear in the settings dialog;
    the GUI LDA tab reads them at start-up to pre-populate its widgets, which
    are then the single control surface for an interactive session.}
  \label{tab:lda-settings}
\end{table}

Turning \code{variable\_projection} off is almost always the wrong choice
(section~\ref{sec:varpro}); it is exposed for diagnostics only.

\subsection{Design}

\code{Settings} is a dataclass and \textbf{the dataclass field is the single
source of truth}. \code{\_\_post\_init\_\_} (which coerces enums),
\code{to\_dict}, \code{from\_dict} and \code{update\_settings} all derive from
\code{dataclasses.fields()}, so adding a setting needs only the field. Enums are
\code{enum.StrEnum}: \code{str(member)} is the value, so they serialise and
compare as plain strings while staying type-checked.

\code{Settings.field\_specs()} drives the GUI settings panel: \code{label},
\code{kind} (\code{choice} / \code{float} / \code{int} / \code{bool} /
\code{text}), optional \code{choices} / \code{min} / \code{max} /
\code{step} / \code{tooltip}, and \code{tab}. The method is \emph{derived}
rather than hand-written: every dataclass field gets an entry automatically,
with the widget kind read from its annotation, the page from
\code{field\_sections()} and the tooltip from \code{field\_comments()}.
\emph{A setting without an entry there does not appear in the panel} ---
which is the intended way to hide an internal knob, and the usual explanation
for a new setting not showing up. A small override table refines labels and
ranges, and can hide a deprecated field.

\subsection{Sections of the file}

\code{settings.toml} is grouped by concern, with a comment above every key:
\code{[fit]} (defaults for a new fit), \code{[solver]} (optimiser limits,
weighting, reporting), \code{[irf]}, \code{[lda]} (LDA defaults for headless
and scripted use), \code{[gui]} and \code{[paths]}. The mapping from field to
section is \code{Settings.field\_sections()}, and it is exhaustive: a field
with no section raises on save rather than being written where the loader would
not find it.

\paragraph{The \code{[lda]} block and the two control surfaces.}
LDA has two ways to set its parameters:
\begin{itemize}
  \item \textbf{Interactive}: the GUI LDA tab --- dropdowns, spinboxes and
    checkboxes. These are the canonical control surface during a session.
  \item \textbf{Headless / scripted}: \code{settings.toml} under \code{[lda]}.
    Keyword arguments to \code{solve\_lda} override everything for a single
    call.
\end{itemize}
\noindent
The \code{[lda]} block does \emph{not} appear in the settings dialog, because
that would create two places to set the same parameters. Instead, the GUI reads
\code{lda\_*} fields from the active \code{Settings} once at start-up and
loads them into the tab's widgets (\code{LDATabMixin.\_init\_lda\_tab\_from\_settings}),
so editing \code{settings.toml} changes the defaults the GUI sees the next time
it opens.

\noindent
The same precedence applies for scripted use:

\begin{lstlisting}
# defaults from settings.toml
result = pr.fit.lda.solve_lda(data)

# or override explicitly --- keyword args always win
result = pr.fit.lda.solve_lda(
    data,
    n_taus=200,
    alpha="auto",
    alpha_method="gcv",
    non_negative=True,
)
\end{lstlisting}

A flat file --- the older layout --- is still read, so an existing
configuration is never orphaned.

\subsection{The configuration file}

Settings are read from and written to \code{settings.toml} with \code{tomlkit},
so user comments survive a save. The file is looked up in the working
directory first, then at the project root.

\begin{lstlisting}
import pyrate_ta as pr

pr.load_settings("settings.toml")
pr.update_settings(n_components=3, irf_mode="Clip", t_min=0.3)
pr.save_settings("settings.toml")
\end{lstlisting}

\code{update\_settings} rejects unknown names rather than storing them, so a
typo fails immediately instead of silently doing nothing.

\subsection{The settings editor}

\code{uv run pyrate-ta-settings} opens a standalone editor whose pages and fields
are generated from \code{field\_specs()}; \emph{Save permanently} writes them
back to \code{settings.toml}. A value the settings layer rejects is logged and
reverted in the widget rather than silently dropped. \code{settings.py} itself
never imports Qt.

The same editor is reachable from the main window under
\emph{View $\rightarrow$ Settings...}, in a dialog with two halves:
\emph{Analysis (PyRATE-TA)} and \emph{Aesthetics (PyMORGAN)}. Both panels edit
their own package's settings live --- the embedded plots re-render as values
change --- and \emph{Save permanently} writes each half back to \emph{its own}
\code{settings.toml}. PyMORGAN's file is located from the installed package
rather than from the working directory: when PyRATE-TA is the host, the
\code{settings.toml} in the working directory is PyRATE-TA's, and writing
PyMORGAN's fields into it would overwrite the wrong file.

\subsection{Both files are loaded}

The GUI loads PyRATE-TA's \code{settings.toml} \emph{and} PyMORGAN's. Without the
second, PyMORGAN's settings stayed at their dataclass defaults, so a profile or
a \code{font\_scale} chosen in PyMORGAN was ignored here and the same data was
drawn with a different font stack --- which is what makes a PyRATE-TA figure look
unlike a PyMORGAN one while both claim the same style. PyMORGAN's file is
located from the installed package, never from the working directory, for the
reason given above.

Fonts themselves are PyMORGAN's (\code{Helvetica, TeX Gyre Heros, Arial}, from
its \code{.mplstyle}) and PyRATE-TA applies them unchanged, but the
\emph{installation} is per environment: a PyRATE-TA virtual environment that has
never run \code{pymorgan-install-fonts} falls back to DejaVu Sans. The GUI
therefore checks at start-up and says so, with the command to run.

\subsection{How results are quoted}

Two settings in \code{[plots]} govern every lifetime PyRATE-TA prints:

\begin{description}
  \item[\code{show\_uncertainties}] Whether the $\pm$ appears at all, in the
    plot legends and the logged fit summary.
  \item[\code{round\_uncertainties}] Round a value \emph{together with} its
    uncertainty: the uncertainty to one significant figure and the value to
    \emph{that same decimal place}, so $325.3 \pm 2.5$ is quoted as
    $325 \pm 3$, $297.5 \pm 1.38$ as $298 \pm 1$ and $8.036 \pm 0.0234$ as
    $8.04 \pm 0.02$. This keeps a fitted constant from being presented with
    more precision than its own uncertainty supports. Halves round \emph{up}
    ($2.5 \rightarrow 3$): an uncertainty should never be rounded down.
\end{description}

Both are applied at the point of display only; the result object keeps the full
precision, so a saved session and the copied lifetimes are unaffected.

\section{Extending PyRATE-TA}\label{sec:extending}

\subsection{Where new code goes}

\begin{description}
  \item[\code{models/}] Kinetic model definitions. A model owns its parameter
    vector layout, its bounds, and the mapping from parameters to \Cmat{}.
    Models are pure NumPy and must be importable without Qt.
  \item[\code{fit/}] The solvers, one module per engine, plus the shared driver
    that handles parameter packing, fixed/free masks and convergence reporting.
    \code{lda.py} is the regularised LDA solver;
    \code{groundstate.py} converts difference spectra to absolute ones;
    \code{constrained.py} provides the sign-mask amplitude solver.
  \item[\code{results/}] Result dataclasses and their serialisation, carrying
    the reproducibility payload of section~\ref{sec:repro}.
    \code{lda\_result.py} is the companion to \code{fit/lda.py}.
  \item[\code{plot/}] Only what PyMORGAN cannot express: residual matrices,
    concentration profiles $\mathbf{c}(t)$, \Kmat{} / compartment diagrams,
    lifetime-density maps, fit-monitor bars, kinetic traces with residuals, and
    spectra with residuals. Everything else is delegated upstream.
  \item[\code{io/}] Import and export of fit sessions --- not raw spectroscopy
    formats, which are PyMORGAN's.
  \item[\code{gui/}] The PyQt6 application (chapter~\ref{sec:gui}).
\end{description}

\subsection{Adding a model}

A model supplies a concentration matrix and the metadata the fit driver and the
result object need: the parameter names, their bounds, and the \code{modelType}
string that will reach \code{plot\_species\_spectra}. Build \Kmat{} explicitly
and reuse the shared propagator (section~\ref{sec:kmat}) rather than deriving a
closed-form solution per scheme --- one propagator is easier to test, and the
degeneracy guard then applies automatically.

Every new model needs a synthetic-recovery test: generate data from known
parameters, fit it, assert the parameters come back within tolerance.

\subsection{Adding an export to the public API}

The public API is lazy (PEP~562). A new export needs an entry in
\code{\_EXPORTS} (name $\rightarrow$ module) \emph{and} in \code{\_\_all\_\_},
and should be added to the \code{TYPE\_CHECKING} block as well so that IDEs and
type checkers still resolve it. Heavy third-party imports used by a single
function (\code{scipy.optimise}, \code{scipy.linalg}, \code{h5py}) belong
\emph{inside} that function.

\subsection{Adding a GUI panel}

Declare the widgets in \code{main\_window.ui} with Qt Designer
(\code{uv run pyrate-ta-edit-gui}), then bind them from a mixin in
\code{gui/tabs/}. Do not build widgets in Python: the layout must stay editable
in Designer. Long operations are wrapped with \code{busy\_guard} so a second
fit cannot start inside the first, and progress is reported through a callback
the GUI adapts, keeping the engine headless.

\subsection{Adding a plot}

New analysis-only figures go in \code{plot/} as a new module, exported through
\code{plot/\_\_init\_\_.py}. Every plotter reads aesthetics from
\code{pymorgan.get\_settings()} so the figures match PyMORGAN's. The module
must not import Qt; the GUI imports the plotter, not the other way around.

\subsection{Conventions}

\begin{description}
  \item[Logging, not \code{print}.] Library code uses
    \code{logger = get\_logger(\_\_name\_\_)}. \code{logger.info} is for
    messages the user should see --- fit summaries, recovered lifetimes,
    convergence status, method citations; \code{logger.debug(\dots,
    exc\_info=True)} goes inside an \code{except} block whose failure is
    recoverable, so that nothing fails silently. \code{print} is reserved for
    console entry points and explicit \code{*\_console} helpers. To debug a
    swallowed failure:
    \code{logging.getLogger("pyrate").setLevel(logging.DEBUG)}.
  \item[Method citations.] A routine implementing a published method --- a
    global-analysis formalism, a target-analysis convention, a regularisation
    scheme --- logs its reference through \code{logger.info} when it finishes.
    An inherited habit from PyMORGAN; keep it.
  \item[No placeholders.] Write fully functional code or raise
    \code{NotImplementedError}. Deprecated code moves to \code{old\_code/}
    (gitignored, and excluded from lint and packaging) rather than being
    deleted.
  \item[Matplotlib.] Use standard mathtext; never enable
    \code{text.usetex=True}.
  \item[Style.] Target Python 3.11, 100-character lines, double quotes:
    \code{uvx ruff@latest format src tests} and
    \code{uvx ruff@latest check src tests scripts examples}.
\end{description}


\section{License}\label{sec:license}
PyRATE-TA is free software released under the \textbf{GNU Affero General Public License version~3.0} (\textsf{AGPLv3}).
You are free to run, inspect, modify and redistribute this software under the terms of the license. Any derivative works or modifications made available over a network or server must also make their complete corresponding source code available under the same license terms.

\end{document}
